\section{Código}
Para esse exercício foi se utilizado programação em linguagem python com integração com Github \textbf{(em anexo no cabeçalho)}.

Para esse exercício foi utilizada a linguagem Python com integração com o GitHub \textbf{(em anexo no cabeçalho)}.

\begin{lstlisting}[caption={Métodos de Newton, Newton Modificado e Newton Discreto}, label={lst:newtons}]
import numpy as np
def newton(x0, y0, Jac, F, itmax=20, tol=1e-6):
    x = x0
    y = y0
    erro = float('inf')
    it = 0

    while erro > tol and it < itmax:
        jac = Jac(x, y)
        if np.linalg.norm(jac) == 0:
            raise np.linalg.LinAlgError("Jacobiana singular no ponto inicial")

        s = np.linalg.solve(jac, -F(x, y))
        x = x + s[0]
        y = y + s[1]
        erro = np.linalg.norm(s)
        it += 1

    return x, y, it, erro


def newton_modificado(x0, y0, Jac, F, itmax=20, tol=1e-6):
    x = x0
    y = y0
    erro = float('inf')
    it = 0

    jac = Jac(x, y)
    while erro > tol and it < itmax:
        s = np.linalg.solve(jac, -F(x, y))
        x = x + s[0]
        y = y + s[1]
        erro = np.linalg.norm(s)
        it += 1

    return x, y, it, erro

# Newton discreto utiliza o método de Newton normal com a jacobiana aproximada
def newton_discreto(x0, y0, Jac_aprox, F, itmax=20, tol=1e-6):
    return newton(x0, y0, Jac_aprox, F, itmax=itmax, tol=tol)
\end{lstlisting}
Declaramos uso de IA (github copilot, gemini, chatGPT e Claude) - para ajuda no código python, LaTeX e dúvidas

Para realizar essa atividade, foi realizado inicialmente a revisão do conteúdo e conseguimos utilizar o método de newton para o primeiro sistema. Após isso, utilizamos IA para nos dar uma visão de como seria a imagem para o primeiro sistema.

Organizamos o código e dividimos entre as funções que calculam os métodos e separamos entre os dois sistemas. Nessa parte utilizamos apenas para conseguir nos ajudar com o matplotlib.

Por fim, estávamos tendo problemas para mostrar no gráfico as raízes do sistema B, foi utilziado para investigar o que estava de errado. Foi utilizado uma tolerância geométrica maior para poder distiguir as 5 regiões das raízes do sistema.